% !TEX root = varianttester.tex


\usetypescriptfile[somerville]
\setuppagenumbering[location=]
%\definepapersize[widepage][width=34cm,height=230cm]
%\setuppapersize[widepage]
\define[1]\MarginStyle{\switchtobodyfont[dejavu,6pt] #1}


\def\opticalsizesall{needlepoint,hairline,titling,display,headline,trumpet,subheading,regular,caption,miniature,miniscule,micro}
\def\weightsall{semilight, none, medium, semibold,black}
\def\widthssall{semicondensed,none,semiexpanded}


%
%
% SETAXES.  Here we set which fonts we want to see
%
%

\def\opticalsizelist{regular,caption}
%\let\opticalsizelist\opticalsizesall

%\def\weightlist{black}
\def\weightlist{none,bold} 
%\let\weightlist\weightsall

\def\widthlist{none}
%\def\widthlist{semiexpanded}
%\def\widthlist{semicondensed}
%\def\widthlist{semiexpanded}

%
%
%
%
%


\def\opticalsizevalue#1{\csname opticalsizevalue:#1\endcsname}

\expandafter\def\csname opticalsizevalue:micro\endcsname{5}
\expandafter\def\csname opticalsizevalue:mico\endcsname{5}
\expandafter\def\csname opticalsizevalue:miniscule\endcsname{6}
\expandafter\def\csname opticalsizevalue:miniature\endcsname{7}
\expandafter\def\csname opticalsizevalue:minature\endcsname{7}
\expandafter\def\csname opticalsizevalue:caption\endcsname{8}
\expandafter\def\csname opticalsizevalue:regular\endcsname{12}
\expandafter\def\csname opticalsizevalue:subheading\endcsname{16}
\expandafter\def\csname opticalsizevalue:trumpet\endcsname{21}
\expandafter\def\csname opticalsizevalue:headline\endcsname{32}
\expandafter\def\csname opticalsizevalue:display\endcsname{41}
\expandafter\def\csname opticalsizevalue:titling\endcsname{48}
\expandafter\def\csname opticalsizevalue:hairline\endcsname{72}
\expandafter\def\csname opticalsizevalue:needlepoint\endcsname{96}

\def\opticalsizenumber{\opticalsizevalue{\opticalsize}}

\setupmargindata
  [inmargin]
  [style=\MarginStyle]





\starttext






\startbuffer[testknuth]
\m{\showglyphs \sum A + \sum B}

\dm{\showglyphs \sum A + \sum B}

\input{knuth}
\startformula
\frac{1}{2\pi i}\oint[bottom=C]{\frac{u(z)dz}{w-z}}= u(w)
\stopformula
\stopbuffer

\startbuffer[test]

\m{
       \dostepwiserecurse{0}{\mathvariantcount "222B}{1} {
       \char\mathvariantcode "222B #1 
       }
 }

\doloopoverlist{.2cm,.8cm,1cm,4cm}{
\dm{\iint[size=\recursestring,bottom=0,top=1,method=horizontal]{u(x) dx} + \int[size=\recursestring,bottom=0,top=1,method=horizontal]{u(x) dx}  }
}
\stopbuffer

\startbuffer[testradical]

   \m{
      \dostepwiserecurse{0}{\mathvariantcount "2A00}{1} {
      \char\mathvariantcode "2A00 #1 
      }
  }

   Here is some text $\sqrt{2} + \sqrt{x}$ so that I can see how text looks like before I play with mathematics
   
   VVV WWWW \m{\odot} \m{\rm \sqrt{2}\sqrt{M}\sqrt{h}\sqrt{L}\sqrt{L^2}}

   \startformula
   \frac{x}{y} \bigodot \frac{1+2+3}{\pi} \bigodot\bigodot \frac{A+b+E+D}{e+f+g+h}
   \stopformula

\m{\sqrt[2]{x} \sqrt[n]{x+M}}

\dm{\sqrt[2]{x} \sqrt[n]{x+M}}

\dm{\sqrt[3]{\frac{A}{B}}}

\dm{\sqrt[3]{\frac{A}{\frac{C}{D}}}}

\startformula

\bigodot[bottom=a,top=b]{{ \frac{\int{\sqrt[n]{ \frac{\bigodot u}{77} }+\sum{ \frac{\sum_q^q u}{\sum v} }}}{\bigodot{ \frac{1}{5} }_{i=1}^n} }}

\sqrt[\math{\displaystyle n^m\hskip-0.9em}]{{ \frac{\int{\sum{ \frac{\sum_q^q u}{77} }+\prod{ \frac{\sum_q^q u}{\prod v} }}}{\bigodot{ \frac{1}{5} }_{i=1}^n} }}
\stopformula

[\hbox to 0pt{\hskip0.5em$A+B$\hss}]

\dm{
\sqrt[\math{n^m}]{x+y}
}
\stopbuffer

\startbuffer[testcirclesimple]
Here is some text UUUOOOUUU 2036

\startformula%\showglyphs
\bigodot_{i=1}^n = \sum_{i=1}^n = \int[bottom=0,top=1]{\frac{f(x)}{g(x)}} dx
\stopformula
  %  \m{
  %     \dostepwiserecurse{0}{\mathvariantcount "2A01}{1} {
  %     \char\mathvariantcode "2A01 #1 
  %     }
  % }

     \m{
      \dostepwiserecurse{0}{\mathvariantcount "2A00}{1} {
      \char\mathvariantcode "2A02 #1 
      }
  }

  %      \m{
  %     \dostepwiserecurse{0}{\mathvariantcount "2A00}{1} {
  %     \char\mathvariantcode "2A02 #1 
  %     }
  % }

% \m{ \char\mathvariantcode "2A01 0 \char\mathvariantcode "2A01 1 \char\mathvariantcode "2A01 2  \char\mathvariantcode "2A01 20}
% \m{A \char\mathvariantcode "2A02 0 \char\mathvariantcode "2A02 2  \char\mathvariantcode "2A02 2 \char\mathvariantcode "2A02 20}

\stopbuffer

\startbuffer[testcircle]
%\startTEXpage

\m{A+B+ C \oplus U \odot A \otimes D \oplus E}



\dm{A+B+ C \oplus U \odot A \otimes D \oplus E}

   \m{
      \dostepwiserecurse{0}{\mathvariantcount "2A01}{1} {
      \char\mathvariantcode "2A02 #1 {\rm U|O}
      }
  }

     \m{
      \dostepwiserecurse{0}{\mathvariantcount "2A00}{1} {
      \char\mathvariantcode "2A00 #1 {\rm U|O}
      }
  }

       \m{
      \dostepwiserecurse{0}{\mathvariantcount "2A09}{1} {
      \char\mathvariantcode "2A09 #1 {\rm U|O}
      }
  }


%  \dm{
%       \dostepwiserecurse{0}{\mathvariantcount "2AFD}{1} {
%      a \char\mathvariantcode "2AFD #1 b,
%       }
%   }

% \dm{
%       \dostepwiserecurse{0}{\mathvariantcount "2AFB}{1} {
%      a \char\mathvariantcode "2AFB #1 b,
%       }
%   }

% \startbuffer[allfrac]
% \getbuffer[tex-frac]
% %\getbuffer[tex-doublefrac]
% \stopbuffer

% \startTEXpage[offset=5pt] 
% \getbuffer[allfrac]
% \par
% \getbuffer[allfrac]

\startformula
\bigtimes = \sum_{i=1}^n
\stopformula

\startformula
\bigtimes{ \frac{\sum u}{\sum v} }_{i=1}^n
\stopformula

\startformula
\bigotimes[bottom=a,top=b]{\sum[bottom=a,top=b]{ \frac{\int[bottom=a,top=b,method=auto]{\bigtimes{ \frac{\sum_q^q u}{77} }+\sum{ \frac{\sum_q^q u}{\sum v} }}}{\sum{ \frac{1}{5} }_{i=1}^n} }}
\stopformula

\startformula
\bigotimes[bottom=a,top=b]{\sum[bottom=a,top=b]{ \frac{\int[bottom=a,top=b,method=auto]{\bigotimes{ \frac{\sum_q^q u}{77} }+\sum{ \frac{\sum_q^q u}{\otimes v} }}}{\sum{ \frac{1}{5} }_{i=1}^n} }}
\stopformula


\startformula
\bigtimes[bottom={i=1},top={i=n}]{{ \frac{\int{\sum{ \frac{\sum_q^q u}{77} }+\prod{ \frac{\sum_q^q u}{\prod v} }}}{\prod{ \frac{1}{5} }_{i=1}^n} }}
\stopformula
%\stopTEXpage
\stopbuffer

\startbuffer[testintegral]
\doloopoverlist{.2cm,.5cm,1cm,2.6cm}{
\m{
\oint[size=\recursestring,bottom=0,top=1]{u}}}
\stopbuffer




% \expandafter\def\csname opticalsizevalue:micro\endcsname{5}
% \expandafter\def\csname opticalsizevalue:miniscule\endcsname{6}
% \expandafter\def\csname opticalsizevalue:miniature\endcsname{7}
% \expandafter\def\csname opticalsizevalue:minature\endcsname{7}
% \expandafter\def\csname opticalsizevalue:caption\endcsname{8}
% \expandafter\def\csname opticalsizevalue:regular\endcsname{12}
% \expandafter\def\csname opticalsizevalue:subheading\endcsname{16}
% \expandafter\def\csname opticalsizevalue:trumpet\endcsname{21}
% \expandafter\def\csname opticalsizevalue:headline\endcsname{32}
% \expandafter\def\csname opticalsizevalue:display\endcsname{41}
% \expandafter\def\csname opticalsizevalue:titling\endcsname{48}
% \expandafter\def\csname opticalsizevalue:hairline\endcsname{72}
% \expandafter\def\csname opticalsizevalue:needlepoint\endcsname{96}



\startbuffer[commutativediagram]

\m{A⇒B⇒C⇒E⇒F⇒G\to H\to I}

\startformula
\lim_{x⇒ 0} f(x)
\stopformula

\startMPcode
%drawarrow (0,0)--(100,0);
label.rt(textext("\m{\rm 0}"), (100,0));
label.rt(textext("\m{\rm A}"), (150,0));
label.rt(textext("\m{\rm B}"), (200,0));
label.rt(textext("\m{\rm D}"), (250,0));
label.rt(textext("\m{\rm C}"), (300,0));
label.rt(textext("\m{\rm 0}"), (350,0));

label.rt(textext("\m{\rm 0}"), (100,50));
label.rt(textext("\m{\rm E}"), (150,50));
label.rt(textext("\m{\rm F}"), (200,50));
label.rt(textext("\m{\rm H}"), (250,50));
label.rt(textext("\m{\rm G}"), (300,50));
label.rt(textext("\m{\rm 0}"), (350,50));

label.rt(textext("\m{\Uhextensible width 1cm \fam 8640}"), (115,0));
label.rt(textext("\m{\Uhextensible width 1cm \fam 8594}"), (165,0));
label.rt(textext("\m{\Uhextensible width 1cm \fam 8594}"), (215,0));
label.rt(textext("\m{\Uhextensible width 1cm \fam 8608}"), (265,0));
label.rt(textext("\m{\Uhextensible width 1cm \fam 8649}"), (315,0));

label.rt(textext("\m{\Uhextensible width 1cm \fam 8611}"), (115,50));
label.rt(textext("\m{\Uhextensible width 1cm \fam 8594}"), (165,50));
label.rt(textext("\m{\Uhextensible width 1cm \fam 8608}"), (215,50));
label.rt(textext("\m{\Uhextensible width 1cm \fam 8594}"), (265,50));
label.rt(textext("\m{\Uhextensible width 1cm \fam 8594}"), (315,50));

label.rt(textext("\m{\Uhextensible width 1cm \fam 8594}") rotated 90, (101,25));
label.rt(textext("\m{\Uhextensible width 1cm \fam 8594}") rotated 90, (151,25));
label.rt(textext("\m{\Uhextensible width 1cm \fam 8608}") rotated 90, (201,25));
label.rt(textext("\m{\Uhextensible width 1cm \fam 8611}") rotated 90, (251,25));
label.rt(textext("\m{\Uhextensible width 1cm \fam 8594}") rotated 90, (303,25));
label.rt(textext("\m{\Uhextensible width 1cm \fam 8594}") rotated 90, (352,25));
\stopMPcode
\stopbuffer




\startbuffer[arrowsidebearings]

\m{→⤅}
%\m{\showglyphs M → M  ← M ↔ M}

\m{\showglyphs M↦M⇥M⤇M}

\stopbuffer

\startbuffer[arrows]

We say that $u$ tends to $0$ as $w$ tends to $0$ (or $u(w)\to 0$ as $w\to 0$) which we also write as 
\startformula
\lim_{w\to 0}u(w) =0 
\stopformula



 \def\MeasureScratchboxInEm{%
  \edef\measuredwidth{%
    \luaexpr{math.floor(1000 * \number\wd\scratchbox
      / \number\fontdimen6\font + .5) / 1000}%
  }%
 }
 \def\ArrowTableCell#1#2{\noexpand\NC H
  \noexpand\setbox\noexpand\scratchbox\noexpand\hbox
    {\noexpand\m{\showglyphs\noexpand\Uhextensible width #2em \noexpand\fam \the\numexpr#1\relax}}%
  \noexpand\MeasureScratchboxInEm
  \noexpand\copy\noexpand\scratchbox
  \noexpand\kern.05em
  {\noexpand\definedfont[Mono at 2pt]%
   \noexpand\measuredwidth em}}
 \starttabulate[|l|l|l|l|l|]
 \doloopoverlist{10521,11066,10576,10571,10570,10574,10502,10523,10498,10599,10603,10590,10582,10594,10602,10598,10586,10578,10500,11069,11060,10522,10513,10517,10516,10503,10524,10499,10601,10605,10591,10583,10600,10604,10596,10587,10579,10501,10497,10518,10520,10519,10496,8594,8592,8596,8602,8603,8695,8696,8697,8698,8699,8622,8606,8608,8610,8611,8612,8614,8676,8677,8633,8636,8637,8640,8641,8651,8652,8644,8646,8649,8647,8658,8656,8660,8653,8654,8655,10229,10230,10231,10232,10233,10234,10235,10236,10237,10238}{
  \xdef\arrowcode{\recursestring}
  \m{ \char\arrowcode\relax } \NC
%  {\definedfont[name:stixtwomathregular at 12pt]\char\arrowcode\relax} \NC
  \doloopoverlist{1,2,4,8}{
   \expanded{\ArrowTableCell{\arrowcode}{\recursestring}}
  }
  \NC\NR
 }
 \stoptabulate


\doloopoverlist{10593,10589,10577,10573,10572,10575,10592,10588,8593,8597,8613,8615,8638,8639,8642,8643,10515,10585,10581,10514,10584,10580}{
  \edef\arrowcode{\recursestring}
  \m{\Uvextensible axis height .5em depth .5em \Umathchar 0 \fam \arrowcode}
}

\doloopoverlist{10593,10589,10577,10573,10572,10575,10592,10588,8593,8597,8613,8615,8638,8639,8642,8643,10515,10585,10581,10514,10584,10580}{
  \edef\arrowcode{\recursestring}
  \m{\Uvextensible axis height 1.5em depth 1.5em \Umathchar 0 \fam \arrowcode}
}

\doloopoverlist{10593,10589,10577,10573,10572,10575,10592,10588,8593,8597,8613,8615,8638,8639,8642,8643,10515,10585,10581,10514,10584,10580}{
  \edef\arrowcode{\recursestring}
  \m{\Uvextensible axis height 2.5em depth 2.5em \Umathchar 0 \fam \arrowcode}
}

\page
\def\arrowcode{10513}

\dostepwiserecurse{100}{290}{5}{%
  \edef\currentwidth{\luaexpr{\recurselevel/100}}%
  \setbox\scratchbox\hbox{%
    \m{\showglyphs\Uhextensible width \currentwidth em \fam \arrowcode}%
  }%
  \MeasureScratchboxInEm
  \framed
    [width=\currentwidth em,
     offset=overlay,
     rulethickness=.1pt,
     framecolor=red]
    {\copy\scratchbox}%
  \kern.05em
  {\definedfont[Mono at 2pt]\measuredwidth em}%
}

\stopbuffer

%
%
% The sample buffers are displayed below
%
%
%




\doloopoverlist\opticalsizelist
{
\edef\opticalsize{\recursestring}
\doloopoverlist\weightlist
{
  \doifelse{\recursestring}{none}
    {\def\weightsuffix{}}
    {\edef\weightsuffix{-\recursestring}}
\doloopoverlist\widthlist
{
  \doifelse{\recursestring}{none}
    {\def\widthsuffix{}}
    {\edef\widthsuffix{-\recursestring}}
\edef\width{\recursestring}
\blank[5pt]
\edef\fontchoice{\opticalsize\widthsuffix\weightsuffix}

\inmargin{{\opticalsize\crlf \widthsuffix\crlf \weightsuffix } @ \opticalsizenumber pt}
{\setupbodyfont[somerville-\fontchoice]
\def\opticalsizepointsize{\opticalsizenumber pt}
%% This next line changes the font size if we want that
%\expanded{\setupbodyfont[\opticalsizepointsize]}
\getbuffer[test]
\getbuffer[testcircle]
\getbuffer[testradical]
%\getbuffer[arrows]
}
}
}
}

\stoptext